% Cote 119: offprints, reports, courses, CV and letters (typed and printed part): archive cote 119, site page 196 (PDF page 197).
% typescript. Draft compared: none.
% AI-transcribed by Claude Opus 5.5 (claude-opus-5-5, adaptive thinking, effort high), Message Batches API, batch r119_full (msgbatch_01DwZNm5Y7YsKBWyAN4sFs8w), 2026-11-13; not yet checked by eye.
\scanpage{119/196}
\uline{ESQUISSE THEMATIQUE DES PRINCIPAUX TRAVAUX MATHEMATIQUES DE A. GROTHENDIECK}

Les numéros entre crochets renvoient, soit à la bibliographie sommaire
jointe à mon Curriculum Vitae (numéros de [1] à [19]), soit au complément à
cette bibliographie placée à la fin du présent rapport (numéros entre [1 bis]
et [20 bis]). Enfin, nous avons joint en dernière page une liste par ordre
alphabétique des auteurs de certains des travaux cités dans le présent
rapport qui ont été directement suscités ou influencés par les travaux de
A. Grothendieck; le renvoi à cette dernière bibliographie se fait par le sigle
[*] derrière le nom de l'auteur cité, comme pour I.M. Gelfand [*].

1. \uline{Analyse Fonctionnelle}([1] à [7], [6 bis])

Mes travaux d'Analyse Fonctionnelle(de 1949 à 1953) ont porté sur-%
tout sur la théorie des espaces vectoriels topologiques. Parmi les nombreuses
notions introduites et étudiées (produits tensoriels topologiques [5,6],
applications nucléaires et applications de Fredholm [5,6,7], applications
intégrales et ses variantes diverses [5,6], applications de puissance p.ème
sommable [5], espaces nucléaires [5], espaces (DF) [4],etc), c'est la notion
d'\uline{espace nucléaire} qui a connu la meilleure fortune : elle a fait jusqu'à
aujourd'hui l'objet de nombreux séminaires et publications. En particulier,un
volume du traité de I. Gelfand [*] sur les ``Fonctions Généralisées'' lui est
consacré. Une des raisons de cette fortune provient sans doute de la théorie
des probabilités, car il s'avère que parmi tous les EVT, c'est dans les espaces
nucléaires que la théorie de la mesure prend la forme la plus simple ( théorème
de Minlos). Les résultats de [6], plus profonds, semblent avoir été moins bien
assimilés par les développements ultérieurs, mais ils apparaissent comme source
d'inspiration dans un certain nombre de travaux délicats assez récents sur des
inégalités diverses liées à la théorie des espaces de Banach, notamment ceux
de Pelczynsky . \quad Signalons également les résultats assez fins de [6] et de
[8 bis] sur les propriétés de décroissance de la suite des valeurs propres de
certains opérateurs dans les espaces de Hilbert et dans les espaces de Banach
généraux.

\uline{Références} : L. Schwartz, J. Dieudonné, I. Gelfand, P. Cartier, J.L. Lions.

\mspage{\hand{196}}
